% Source PDF: source/普通高考/2019/2019天津文.pdf, page 1, question 4.
% PDF embedded-bitmap check: pdfimages -list found no image objects; the flowchart is vector artwork.
% Grid evidence: tmp/redraw_flowchart_2019_tianjin_liberal/source_p1_grid100.jpg;
% comparison original side: tmp/redraw_flowchart_2019_tianjin_liberal/q4_original.png,
% cropped from the no-grid page render.
% Elements: 开始/结束 rounded boxes, process rectangles, decision diamonds,
%          输出 S parallelogram, branch labels, arrows, and loop-back branch.
% Node→edge list (from source PDF):
%   开始→i=1,S=0→i为偶数?;
%   是→j=i/2→S=S+i·2^j→i=i+1;
%   否→S=S+i→i=i+1;
%   i=i+1→i≥4?; 否→回到 i为偶数?; 是→输出S→结束.
% solid segments: all process/decision/output outlines, loop-back and branch lines, arrows.
% dashed segments: none.
% labels: branch 否/是 are offset from decision edges and arrows; formulas stay centered.

\begin{tikzpicture}[baseline=(current bounding box.center), x=1cm, y=1.18cm,
  flowbox/.style={draw, anchor=center, line width=0.45pt, minimum height=0.55cm,
    inner xsep=4pt, inner ysep=2pt, align=center, font=\normalsize,
    execute at begin node=\setlength{\parindent}{0pt}},
  flowarrow/.style={-Stealth, line width=0.45pt},
]
  % The source is a schematic flowchart; coordinates preserve topology and branch order.
  \node[flowbox, rounded corners=0.12cm, minimum width=1.10cm] (start) at (0,0) {开始};
  \node[flowbox, minimum width=2.70cm] (init) at (0,-1.08) {$i=1,S=0$};
  \node[draw, line width=0.45pt, diamond, aspect=2.75, minimum width=5.10cm,
    minimum height=1.28cm, inner xsep=4pt, inner ysep=2pt, anchor=center, align=center, font=\normalsize,
    execute at begin node=\setlength{\parindent}{0pt}] (even) at (0,-2.50)
    {$i$ 为偶数?};
  \node[flowbox, minimum width=1.62cm, minimum height=0.90cm] (jcalc) at (0,-4.05)
    {$j=\dfrac{i}{2}$};
  \node[flowbox, minimum width=3.25cm] (seven) at (0,-5.42)
    {$S=S+i\cdot2^j$};
  \node[flowbox, minimum width=2.40cm] (sodd) at (3.10,-4.05)
    {$S=S+i$};
  \node[flowbox, minimum width=1.95cm] (inc) at (0,-6.80)
    {$i=i+1$};
  \node[draw, line width=0.45pt, diamond, aspect=2.75, minimum width=4.55cm,
    minimum height=1.18cm, inner xsep=4pt, inner ysep=2pt, anchor=center, align=center, font=\normalsize,
    execute at begin node=\setlength{\parindent}{0pt}] (gefour) at (0,-8.25)
    {$i\geq4?$};
  \node[flowbox, minimum width=2.20cm, minimum height=0.58cm,
    trapezium, trapezium left angle=75, trapezium right angle=105] (output) at (0,-9.62)
    {输出 $S$};
  \node[flowbox, rounded corners=0.12cm, minimum width=1.10cm] (finish) at (0,-10.82) {结束};

  \draw[flowarrow] (start) -- (init);
  \draw (init.south) -- (0,-1.60);
  \draw[flowarrow] (0,-1.60) -- (even.north);
  \draw[flowarrow] (even) -- node[font=\normalsize, right, xshift=0.05cm] {是} (jcalc);
  \draw[flowarrow] (jcalc) -- (seven);
  % Keep the branch stem clear of the odd-case box before attaching to it.
  \draw[line width=0.45pt] (even.east) -- (3.10,-2.50) -- (3.10,-3.30);
  \draw[flowarrow] (3.10,-3.30) -- (sodd.north);
  \node[font=\normalsize, anchor=west, inner sep=0.02cm] at (2.55,-2.13) {否};
  % Both branches join immediately above i=i+1; the horizontal arrow points left
  % and the odd-branch arrow terminates at the visible merge point.
  \draw[line width=0.45pt] (seven.south) -- (0,-6.20);
  \draw[flowarrow] (sodd.south) -- (3.10,-6.20) -- (0.30,-6.20) -- (0,-6.20);
  \draw[flowarrow] (0,-6.20) -- (inc.north);
  \draw[flowarrow] (inc) -- (gefour);
  % The loop returns into the trunk above the decision; the main trunk carries
  % the sole downward arrowhead into the decision.
  \draw[line width=0.45pt] (gefour.west) -- (-3.30,-8.25) -- (-3.30,-1.60);
  \draw (-3.30,-1.60) -- (-0.75,-1.60);
  \draw[flowarrow] (-0.75,-1.60) -- (0,-1.60);
  \node[font=\normalsize, anchor=east, inner sep=0.02cm] at (-2.70,-8.00) {否};
  \draw[flowarrow] (gefour) -- node[font=\normalsize, right, xshift=0.06cm] {是} (output);
  \draw[flowarrow] (output) -- (finish);
  % Preserve the source's safe white margin around the outer loop and end boxes.
  % Include the full right-hand odd branch and a small even margin around the end box.
  \path[use as bounding box] (-3.55,0.45) rectangle (4.55,-11.20);
\end{tikzpicture}
